\section{Additional Analysis for Course Submission}

\subsection{Why the result is a sufficient-condition baseline}

The central theorem is intentionally a baseline rather than a complete theory
of MoE specialization. It says that if token directions follow the spherical
model and if the router prototypes induce equal-measure spherical cells, then
unconstrained top-1 routing has uniform expected usage. The value of this
statement is that it makes the hidden assumptions behind a tempting load-balance
interpretation explicit. If a trained router shows balanced load, the theorem
does not by itself certify that the model has found semantic features. Instead,
it gives a list of objects that must be audited before such an interpretation is
safe: the token direction distribution, the prototype-cell geometry, the
finite-token sampling error, and any capacity or auxiliary-loss mechanism that
may have forced the counts.

This distinction is important because the same observed load statistic can have
several causes. A low LoadCV can arise because the representation is close to
spherical and the prototypes are well spread. It can also arise because a
capacity constraint has forced equal counts, because an auxiliary loss has
penalized imbalance, or because a small batch happened to sample nearly equal
counts. Conversely, a high LoadCV can arise from nonspherical token directions,
from clustered prototypes, from finite-token noise, or from a genuine training
dynamic that amplifies an early routing advantage. The controlled model in this
paper isolates the first two geometric terms and treats the other mechanisms as
separate explanations rather than background noise.

\begin{table}[t]
\centering
\caption{Operational role of the four uniformity notions.}
\label{tab:uniformity_audit}
\small
\setlength{\tabcolsep}{3pt}
\resizebox{\linewidth}{!}{%
\begin{tabular}{p{0.25\linewidth}p{0.31\linewidth}p{0.33\linewidth}}
\toprule
Notion & What is measured & What it cannot prove \\
\midrule
Feature-frequency uniformity & Labels or feature families occur with similar
probability. & Hidden directions are spherical or routing will balance. \\
Directional uniformity & Centered, normalized routing states cover the sphere
approximately evenly. & Expert cells have equal measure. \\
Spectral uniformity & The covariance spectrum has no dominant direction. &
The directional distribution is spherical. \\
Usage uniformity & Empirical expert loads are close to $1/E$. & The route is
semantic, useful, or naturally produced. \\
\bottomrule
\end{tabular}
}
\end{table}

Table~\ref{tab:uniformity_audit} summarizes the operational meaning of the four
uniformity notions. The table is not just terminology. It is the paper's main
guard against metric substitution. A count metric is only the final outcome of
the router. It cannot be used as evidence for the token-side representation
unless the assignment mechanism is known and the rival explanations have been
separated.

\subsection{Reading the experiments as falsification checks}

The six experiments should be read as falsification checks against nearby but
weaker explanations. The first experiment tests the positive sufficient
construction: if the token distribution is spherical and the prototypes are
well spread, top-1 usage is nearly uniform. This is the clean case that the
theory predicts. The result does not say that all random or trained prototypes
will be well spread. It says that well-spread cells are enough when the
token-side condition holds.

The second experiment separates the prototype-cell floor from finite-token
sampling noise. This matters because a finite experiment always contains both
population geometry and sampling error. If increasing the number of routed
examples removes the whole imbalance, then the problem was mostly sampling. If
the total error approaches a nonzero floor, then the fixed prototype geometry is
the remaining cause. The random-prototype floor experiment shows the latter
behavior: larger batches reduce the sampling term but do not remove the cell
mass error for a fixed random prototype realization.

The third experiment tests the expert-prototype interpretation. Singular
directions of an expert matrix can be meaningful centers, and a power-iteration
style update can align router rows with those centers. The experiment is
therefore not a negative statement about spectral expert geometry. Its narrower
point is that center recovery and load balance are different properties. If the
recovered centers cluster, then top-1 cells can still be imbalanced even when
alignment is high. This is the reason the paper treats spectral prototypes as a
construction rule for candidate centers, not as a theorem for balanced routing.

The fourth experiment is a counterexample to an overly strong spectral claim.
The four-axis distribution in $\R^2$ has identity covariance up to scale, but
its support is concentrated on four directions. Regular-simplex prototypes have
equal cells for the continuous uniform circle distribution, yet the discrete
four-axis distribution places unequal mass in those cells. The counterexample
shows why covariance isotropy is only a second-order diagnostic. It may detect a
large anisotropy, but it cannot replace a directional support audit.

The fifth experiment is deliberately limited to an elliptical toy family. In
that family, whitening inverts the known linear distortion and recovers a
spherical direction distribution up to rotation. The success of whitening in
this experiment is therefore evidence for Proposition~\ref{prop:whitening}, not
evidence that arbitrary language-model hidden states are generated by a single
elliptical distortion. This boundary prevents a common overinterpretation:
spectral flattening can be a useful repair when its model assumptions hold, but
it is not a universal certificate of routable uniformity.

The sixth experiment treats capacity assignment as a control rather than as the
mechanism of interest. Balanced assignment can make the count vector exactly
uniform by changing the assignment problem. This is useful when the engineering
goal is to keep all experts active in a batch, but it does not explain why an
unconstrained top-1 router would naturally find equal-mass cells. The score
drop in the assignment objective records the price paid to force the count
constraint. Thus the experiment separates count equality from representation
geometry.

\subsection{Relation to the presentation line}

The presentation material prepared earlier for this course studies a broader
training-time question: why sparse MoE routers may concentrate before useful
expert specialization forms. That story includes common-logit bias,
router-hidden-state amplification, and expert-spectrum initialization stress
tests. The present manuscript is narrower. It does not use those real-text
training observations as proof of the main theorem. Instead, it provides the
static geometric baseline that the broader project needs before interpreting
load balance or route consistency.

There is therefore no logical contradiction between the slide story and this
paper if their claim strengths are kept separate. The slides ask which geometry
is amplified first during early training. This paper asks what must be true in a
static controlled router before uniform usage can be attributed to
representation geometry. The slide evidence can motivate the need for a
geometry audit, but it should not be used to claim that the sufficient condition
holds in real text. Conversely, the theorem in this paper can clarify what
would need to be checked in the training setting, but it does not solve
training-time preservation, expert utility, or semantic feature identity.

This separation also explains why the paper avoids stronger language such as
``experts specialize uniformly'' or ``spectral initialization fixes routing.''
Uniform load is a count-level observation. Specialization requires a stronger
link between routed tokens and useful expert functions. Training preservation
requires a margin argument that controls how router rows, token states, and
expert parameters move under optimization. Neither property follows from the
static equal-measure theorem. The static theorem is still useful because it
gives the null model: before claiming specialization from load, one must first
show that the observed load is not merely a consequence of cell geometry,
sampling, or forced assignment.

\subsection{A practical audit pipeline}

The controlled analysis can be used as a practical audit pipeline for future
MoE routing experiments. The pipeline has five stages. Each stage answers one
question and deliberately refuses to answer the next question. This staged
reading is useful because many MoE routing metrics are easy to overinterpret.

The first stage is the population stage. The object is not yet the router, but
the token set that will be routed. One must define which hidden states are
route-relevant, how they are centered, and whether the analysis uses all
positions or a restricted feature population. If the wrong population is used,
then a balanced or imbalanced router statistic may describe an irrelevant
mixture of states. In the controlled experiments, the population is explicit:
the examples are sampled from the sphere, from a fixed random-prototype cell
estimator, from a four-axis distribution, or from an elliptical transformation
of a spherical source. This explicit population definition is why the
experiments can support precise but narrow claims.

The second stage is the directional stage. The object is
$z=(h-\mu)/\|h-\mu\|_2$, not the raw hidden state $h$. The top-1 score compares
inner products, so after centering and scale normalization the angular
distribution is the part that directly controls cell mass. A covariance
spectrum is useful here only as a diagnostic. It can warn that the distribution
has a large dominant direction, and it can suggest whether a whitening repair is
worth trying under an elliptical model. It cannot prove that every angular
region has the right mass. The four-axis counterexample is the smallest
operational demonstration of this gap.

The third stage is the prototype-cell stage. Given prototypes
$u_1,\ldots,u_E$, the relevant object is the partition
$V_1,\ldots,V_E$ of the sphere. Uniform expected load requires equal mass under
the token direction distribution, not merely large pairwise angles or high
singular-vector alignment. Regular-simplex prototypes give a clean positive
case because symmetry forces congruent cells under the spherical distribution.
Random prototypes do not generally give this guarantee in one realization.
Clustered prototypes explicitly violate it. Expert singular directions are
therefore useful only after the induced cell masses are checked.

The fourth stage is the finite-sample stage. The empirical load vector
$\widehat p$ is a noisy estimate of the population cell masses $p$. The
decomposition
\[
\widehat p_e-\frac{1}{E}=
\left(p_e-\frac{1}{E}\right)+(\widehat p_e-p_e)
\]
is the audit rule. If the geometry term is large, increasing the batch size
cannot fix the underlying cell imbalance. If the sampling term is large, more
tokens can reduce the empirical error. This is why the random-prototype floor
experiment reports geometry, sampling, and total error separately instead of
only one final LoadCV.

The fifth stage is the mechanism-control stage. If a capacity rule, auxiliary
loss, or assignment solver forces the count vector, then the count vector no
longer identifies the natural top-1 geometry. Such controls are not wrong; they
are often useful engineering devices. They simply answer a different question:
what score cost is required to enforce equal counts? The balanced-assignment
experiment keeps this control visible so that forced count equality is not
mistaken for a representation-geometric explanation.

\begin{table}[t]
\centering
\caption{Experiment-to-question map. Each experiment removes one nearby
overinterpretation.}
\label{tab:experiment_question_map}
\small
\setlength{\tabcolsep}{3pt}
\resizebox{\linewidth}{!}{%
\begin{tabular}{p{0.23\linewidth}p{0.34\linewidth}p{0.34\linewidth}}
\toprule
Experiment & Question answered & Overclaim blocked \\
\midrule
Symmetric sphere & Can equal cells balance spherical tokens? &
Balanced load is impossible under top-1. \\
Random floor & Does more data remove fixed cell imbalance? &
All imbalance is finite-sample noise. \\
Spectral centers & Does expert-direction alignment imply balance? &
High alignment is a load theorem. \\
Covariance example & Does flat covariance imply uniform routing? &
Second-order isotropy certifies directions. \\
Whitening toy & Can known elliptical distortion be repaired? &
Whitening is universally valid. \\
Balanced assignment & Can count equality be forced? &
Forced counts prove natural routing. \\
\bottomrule
\end{tabular}
}
\end{table}

Table~\ref{tab:experiment_question_map} gives the intended reading order. The
experiments are not six independent attempts to improve a single score. They are
six controls around one interpretation problem. The paper asks when uniform
expert usage is evidence for a natural top-1 geometric mechanism. Each row
removes one explanation that would otherwise make the count metric ambiguous.

\subsection{Implications for future real-model work}

The next step outside this paper is not to assume that real language-model
states satisfy the spherical model. The next step is to measure which part of
the audit pipeline fails. A real model may fail at the population stage because
the routed states mix several unrelated token populations. It may fail at the
directional stage because hidden states contain a common component or a few
dominant prediction directions. It may fail at the prototype-cell stage because
router rows or expert-derived centers cluster. It may fail at the finite-sample
stage if the evaluation batch is too small to separate sampling noise from a
fixed cell floor. It may also pass the count audit while failing the utility
audit because routed tokens are not assigned to experts that perform different
useful computations.

This paper therefore suggests a conservative decision rule for future work.
Before interpreting uniform load as specialization, first check whether the
same conclusion survives after defining the routed population, auditing
directional mass, estimating prototype-cell imbalance, separating finite-sample
noise, and ruling out forced assignment. If the answer is yes, then load
uniformity becomes a more meaningful signal. If the answer is no, then the
failure location is itself informative: it tells whether the obstacle is token
geometry, prototype geometry, sampling, assignment control, or a later
training-time utility question.

The presentation-line experiments fit into this decision rule as downstream
motivation. They suggest that early training can re-amplify a routing bias even
when step-zero initialization is improved. The static theory here does not
explain that dynamic by itself. What it provides is the baseline that the
dynamic theory must exceed. A future training-time claim should therefore add
new objects: margin lower tails, router-row drift, token-state drift, and loss
or utility guards. Without those objects, a dynamic result would again risk
collapsing several explanations into one load number.

\subsection{Reproducibility and source map}

For a course submission, the most important reproducibility requirement is that
the reader can connect every claim to a small source object. This paper keeps
the source objects simple: all experiments are synthetic, all central figures
come from the same controlled variables described in the text, and no claim
depends on a hidden checkpoint or a private dataset. The expected reproduction
path is therefore not a large training run. It is a sequence of deterministic
or seeded simulations that instantiate the mathematical objects in the theory.

\begin{table}[t]
\centering
\caption{Source map for the main evidence objects.}
\label{tab:source_map}
\small
\setlength{\tabcolsep}{3pt}
\resizebox{\linewidth}{!}{%
\begin{tabular}{p{0.26\linewidth}p{0.34\linewidth}p{0.31\linewidth}}
\toprule
Evidence object & Controlled variables & Reader check \\
\midrule
Prototype geometry & $d$, $E$, prototype spread, spherical tokens &
Symmetric cells have near-zero LoadCV. \\
Geometry floor & Random prototype realization, reference cell mass, $N$ &
Sampling error falls while geometry error remains. \\
Spectral centers & Signal directions $v_e$, noise level, MPI steps &
Alignment and load balance are separated. \\
Covariance example & Four-axis support and three simplex prototypes &
Covariance is flat but cell mass is unequal. \\
Whitening toy & Elliptical scale $\alpha$ and whitening transform &
Repair works only under the stated model. \\
Balanced assignment & Batch size, scores, capacity constraint &
Counts are forced at a measurable score cost. \\
\bottomrule
\end{tabular}
}
\end{table}

Table~\ref{tab:source_map} is also a guard against hidden overclaiming. A reader
should be able to ask, for each row, which variable was changed and which rival
explanation was blocked. If a figure cannot answer that question, it should be
treated as a diagnostic visualization rather than decision evidence. This is
why the paper uses a small number of figures and keeps the captions tied to the
claim boundary.

The source map also clarifies what is intentionally absent. There is no
language-model loss curve because language-model utility is outside the static
claim. There is no semantic-label purity score because semantic expert identity
is outside the controlled feature-routing model. There is no training-time
preservation metric because optimization drift is a later question. These
missing measurements are not oversights; they are excluded so that the paper can
make a clean statement about the static geometry before adding the harder
dynamic variables.

If this work is extended to real checkpoints, the source map should become more
strict rather than less strict. Each added real-model figure should define the
token population, checkpoint, layer, routing input, centering or projection
operator, router rows, expert matrices, and batch size. It should also report
the false-positive guard that prevents a load or route-consistency metric from
being mistaken for functional specialization. The controlled paper therefore
sets the reporting standard for future work: every metric must say which
uncertainty it reduces and which uncertainty it leaves open.

\subsection{Claim boundary and falsifiers}

The allowed claim after the evidence in this paper is that uniform usage under
unconstrained top-1 routing has a clean sufficient geometric explanation in a
controlled setting. Token directions must match the spherical model, and expert
prototypes must induce equal-measure cells. Spectral diagnostics can help locate
anisotropy and construct candidate centers, but they are not interchangeable
with the directional and cell-mass conditions. Capacity assignment can force the
target count vector, but it changes the mechanism being studied.

Several observations would falsify or narrow this account. If symmetric
regular-simplex prototypes failed to balance a high-dimensional spherical
source beyond finite sampling noise, then the equal-measure implementation
would be wrong. If increasing $N$ removed the entire random-prototype floor,
then the decomposition into geometry error and sampling error would be
misapplied. If covariance identity always implied uniform routing for
noncontinuous directional distributions, then the counterexample would fail.
If whitening did not repair the controlled elliptical family at the population
level, then the whitening proposition or its implementation would be wrong.
These are concrete failure modes because each one attacks a specific term in
the model.

Other observations would not falsify the static theorem but would limit its use
in real training. For example, a trained language model may have strong
anisotropy that is not elliptical, or the route-relevant population may be only
a small subset of all hidden states. A model may preserve balanced load while
collapsing distinct features into the same expert, or it may have high route
consistency without useful expert functions. These cases do not contradict the
controlled theorem. They show why the theorem is only the first layer of the
larger specialization audit.
